\documentclass[../../main.tex]{subfiles}
\begin{document}


{\bf [解]}
与えられた関数
\begin{align}
 f(x)=x^3+ax^2+(b-a-1)x \label{eq:1}
\end{align}
を考える．

\paragraph{(1)}
$f(x)$の挙動を調べるために一階微分を計算すると
\begin{align}
 f'(x)=3x^2+2ax+(b-a-1) = 3\left(x+\dfrac{a}{3}\right)^2 + (b-a-1) -\dfrac{a^2}{3} \label{eq:2}
\end{align}
である．題意を満たす$(a,b)$はこれが$x\ge0$で0以上であるような時である．
$f'(x)$は二次関数だから軸の位置$x_0=-\dfrac{1}{3}a$の大きさで場合分けする．


\subparagraph{$1^\circ$ $-\dfrac{1}{3}a\le0\ \therefore\ a\ge0$の時}

この時は$x \ge 0$で$f'$は単調増加だから，$f'(x) \ge 0$と成る条件は
\begin{align}
 f'(0)\ge0\iff b\ge a+1 \label{eq:3}
\end{align}
である．

\subparagraph{$2^\circ$ $0\le-\dfrac{1}{3}a\ \therefore\ a\le0$の時}

この時は$0 \le x_0$で$f'$は最小値をとるから，これが非負ならば$f'(x) \ge 0$となるから条件は
\begin{align}
 f'\left(-\dfrac{a}{3}\right) \ge 0\iff b\ge\dfrac{1}{3}a^2+a+1 \label{eq:4}
\end{align}
である．
\casesend

以上\cref{eq:3,eq:4}が求める条件でまとめると
\begin{align}
\begin{cases}\displaystyle
a\le0\text{の時}\ b\ge\dfrac{1}{3}a^2+a+1 \\
a\ge0\text{の時}\ b\ge a+1
\end{cases} \label{eq:5}
\end{align}
となる．図示すると下図斜線部である．

\begin{figure}[htb]
\begin{center}
\begin{tikzpicture}[scale=1.2, >=stealth]

  % --- 1. 座標軸 ---
  \draw[->, thick] (-4.2,0) -- (2.2,0) node[right] {$a$};
  \draw[->, thick] (0,-0.5) -- (0,3.8) node[above] {$b$};
  \node[below left] at (0,0) {$O$};

  % --- 2. 求める領域 G の斜線ハッチング (pattern=north east lines) ---
  % a <= 0: b >= a^2/3 + a + 1
  % a >= 0: b >= a + 1
  \begin{scope}
    \clip (-4.2,0) rectangle (2.2,3.6);
    \fill[pattern=north east lines, pattern color=blue!60]
      (-4.2,3.6) -- plot[domain=-4.2:0, samples=50] (\x, {\x*\x/3 + \x + 1})
      -- plot[domain=0:2.2, samples=50] (\x, {\x + 1})
      -- (2.2,3.6) -- cycle;
  \end{scope}

  % --- 3. 境界線の描画 ---
  % (i) a <= 0 側の放物線 b = a^2/3 + a + 1
  \draw[ultra thick, blue!80!black, domain=-4.0:0, samples=100] 
    plot (\x, {\x*\x/3 + \x + 1}) node[above left, font=\small] {$b = \frac{1}{3}a^2+a+1$};

  \draw[blue!80!black, domain=0:2.0, samples=100] 
    plot (\x, {\x*\x/3 + \x + 1}); 

  % (ii) a >= 0 側の直線 b = a + 1
  \draw[ultra thick, blue!80!black, domain=0:2.2, samples=50] 
    plot (\x, {\x + 1}) node[above left, font=\small] {$b = a+1$};

  \draw[blue!80!black, domain=-2.0:0, samples=50] 
    plot (\x, {\x + 1});


  % --- 4. 特徴的な点と目盛りのプロット ---
  % 放物線の頂点 a = -3/2 = -1.5, b = 1/4 = 0.25
  \draw[dashed, gray] (-1.5,0) node[below, font=\small] {$-\frac{3}{2}$} -- (-1.5, 0.25) -- (0, 0.25) node[right, font=\small] {$\frac{1}{4}$};
  \fill[black] (-1.5, 0.25) circle (1.2pt);

  % a 軸との交点 a = -3
  \fill[black] (-3,0) circle (1.2pt) node[below, font=\small] {$-3$};

  % b 軸片 (0, 1)
  \fill[black] (0,1) circle (1.5pt) node[left, font=\small] {$1$};

  % 領域名 G
  \node[blue!80!black, font=\large] at (-1.5, 2.2) {$G$};

\end{tikzpicture}
\end{center}
\caption{条件を満たす $(a,b)$ の範囲 $G$（斜線部・境界含む）}
\end{figure}

\paragraph{(2)} 

$(a,b)$が$G$をうごく時(1)より$0\le x$で$f(x)$は単調増加かつ$f(1)=b$だからグラフの概形は下図である．
求める面積$S$は図の斜線部である．

\begin{figure}[htb]
\begin{center}
\begin{tikzpicture}[scale=1.6, >=stealth]

  % --- 1. 座標定義 ---
  \coordinate (O) at (0,0);
  \coordinate (A) at (1,0);
  \coordinate (B) at (1,0.85);
  \coordinate (C) at (0,0.85);

  % --- 2. 座標軸 ---
  \draw[->, thick] (-0.3,0) -- (2.2,0) node[right] {$x$};
  \draw[->, thick] (0,-0.3) -- (0,2.0) node[above] {$y$};
  \node[below left] at (O) {$O(0,0)$};

  % --- 3. 斜線ハッチング領域 (pattern=north east lines) ---
  % y = f(x) の上側かつ x=0, y=b で囲まれた領域
  % ※例として a = -0.5, b = 1.5 付近の単調増加曲線を使用
  \begin{scope}
    \clip (0,0) -- plot[domain=0:1, samples=50] (\x, {0.5*\x*\x*\x - 0.25*\x*\x + 0.6*\x})
          -- (0,0.85) -- cycle;
    \fill[pattern=north east lines, pattern color=blue!60] 
      (0,0) rectangle (1,0.85);
  \end{scope}

  % --- 4. 曲線 y = f(x) の描画 ---
  \draw[ultra thick, blue!80!black, domain=0:1.3, samples=100] 
    plot (\x, {0.5*\x*\x*\x - 0.25*\x*\x + 0.6*\x}) node[right] {$y = f(x)$};

  % --- 5. 長方形 OABC の枠線・破線 ---
  % 上辺 CB (実線または破線)
  \draw[thick, gray] (C) -- (B);
  % 右辺 AB (破線)
  \draw[dashed, gray] (A) -- (B);

  % --- 6. 各頂点のプロットとラベル ---
  \fill[black] (A) circle (1.2pt) node[below] {$A(1,0)$};
  \fill[black] (B) circle (1.5pt) node[above right] {$B(1,b)$};
  \fill[black] (C) circle (1.2pt) node[left] {$C(0,b)$};

\end{tikzpicture}
\end{center}
\caption{曲線 $y=f(x)$ と求める面積$S$の領域（斜線部）}
\end{figure}

新しく点を$O(0,0), A(1,0), B(1,b), C(0,b)$とおくと面積$S$は
\begin{align}
\begin{split}
S
 &= \int_{0}^{b}f^{-1}(y) dy \\
 &=\text{四角形} OABC-\int_0^1f(x)dx \\
 &=b-\left(\dfrac{1}{4}+\dfrac{1}{3}a+\dfrac{1}{2}(b-a-1)\right) \\
 &=\dfrac{1}{2}b+\dfrac{1}{6}a+\dfrac{1}{4} \quad\cdots\text{①}
\end{split} \label{eq:6}
\end{align}
である．この最小値を求めるため，\cref{eq:5}から$a$の値で場合分けする．

\subparagraph{$a\le0$の時}

この時は\cref{eq:5}から$b\ge\dfrac{1}{3}a^2+a+1$だから\cref{eq:6}を下から評価して
\begin{align}
\begin{split}
 S
 &\ge \dfrac{1}{6}a+\dfrac{1}{2}\left(\dfrac{1}{3}a^2+a+1\right)+\dfrac{1}{4} \\
 &=\dfrac{1}{6}a^2+\dfrac{2}{3}a+\dfrac{3}{4} \\
 &=\dfrac{1}{6}(a+2)^2+\dfrac{1}{12} \\
 &\ge\dfrac{1}{12} \quad(\text{等号成立は}a=-2)
\end{split} \label{eq:7}
\end{align}
である．

\subparagraph{$a\ge0$の時}

この時は\cref{eq:5}から $b\ge a+1$ だから\cref{eq:6}を下から評価して
\begin{align}
\begin{split}
S
&\ge \dfrac{1}{6}a+\dfrac{1}{2}(a+1)+\dfrac{1}{4} \\
&=\dfrac{2}{3}a+\dfrac{3}{4} \\
&\ge\dfrac{3}{4} \quad(\text{等号成立は}a=0)
\end{split} \label{eq:8}
\end{align}
である．
\casesend

以上\cref{eq:7,eq:8}から求める最小値は$a=-2, b=\dfrac{4}{3}$の時の
\begin{align}
 S=\dfrac{1}{12}
\end{align}
である．

\newpage

\end{document}
